
\begin{table}[t]
\centering
\caption{Dataset Composition: Four-Stage Research Pipeline}
\label{tab:dataset_composition}
\small
\begin{tabular}{@{}lrll@{}}
\toprule
\textbf{Dataset} & \textbf{Sample Size (N)} & \textbf{Reward Labels?} & \textbf{Role in Study} \\
\midrule
RouteLLM Battles & 80,000 & Yes & Training: Initializing the Warmup Prior. \\
LMSYS Dev Set & 1,121 & Yes & Validation: Performance tuning \& Pareto optimization. \\
LMSYS Holdout Set & 750 & Yes & Testing: Final unseen performance metrics (Table 2). \\
LMSYS Chat-1M & 1,000,000 & No & Scaling: Manifold validation \& ROI estimation. \\
\bottomrule
\end{tabular}
\vspace{0.5em}

\begin{tablenotes}
\small
\item \textbf{Four-Stage Research Pipeline:}
\item \textit{Stage 1 -- Training (80k RouteLLM Battles)}: LMSYS Arena battles from RouteLLM dataset \cite{routellm2024} with reward labels. Establishes the foundation for PCA training (384→32 dims) and LinUCB warmup priors (covariance matrix $\mathbf{A} \in \mathbb{R}^{33 \times 33}$ and belief vector $\mathbf{b} \in \mathbb{R}^{33}$). Provides sufficient scale to learn robust semantic representations.
\item \textit{Stage 2 -- Validation (1,121 LMSYS Dev)}: Development set with ground-truth reward labels for performance tuning, Pareto optimization, and Corralling calibration. Used for Figure~\ref{fig:corralling_semantic} (Right) expert weight evolution.
\item \textit{Stage 3 -- Testing (750 LMSYS Holdout)}: Held-out test set with ground-truth reward labels for computing final Regret and AUPR metrics (Table~\ref{tab:performance-gap}). Completely disjoint from training and validation data, ensuring unbiased performance measurement.
\item \textit{Stage 4 -- Scaling (1M LMSYS Chat)}: Large-scale unlabeled dataset proving semantic manifold stability at production scale. Used for Figure~\ref{fig:corralling_semantic} (Left) to validate that PCA-learned representations generalize beyond labeled training data to real-world prompt distributions. Enables ROI estimation and demonstrates the 94.2\% Easy Cluster dominance.
\item \textbf{LinUCB Priors:} Warmup data initializes the contextual bandit with covariance matrix $\mathbf{A}$ (capturing feature correlations) and belief vector $\mathbf{b}$ (encoding reward expectations) for each model. Context dimension: 33 (32 PCA components + 1 bias term).
\item \textbf{Quality Assurance:} All labeled datasets (Stages 1-3) are completely disjoint, preventing data leakage. The 1M dataset is used exclusively for visualization and scaling analysis, not for training or evaluation.
\end{tablenotes}
\end{table}
